% ============================================================================
% DM-2026-010: PCFPlan Increment Editor Enhancements
%   Decision 1 — Increment display-window selection (REQ-011a / issue #32)
%   Decision 2 — Payload-config copy collision behavior (REQ-060 / issue #33)
% ============================================================================
\newcommand{\UniqueID}{DM-2026-010}
\newcommand{\DocumentDate}{July 23, 2026}
\newcommand{\AuthorName}{Bruce Dombrowski}
\newcommand{\AuthorTitle}{Systems Engineer, Payload Software Integration}
\newcommand{\ToField}{PCFPlan Stakeholders, ISS Payload Integration}
\newcommand{\SubjectField}{PCFPlan Increment Editor: Display-Window Selection and Payload-Config Copy}
\newcommand{\OPRField}{Payload Software Integration (PSI)}

\newcommand{\dmContent}{%
\begin{enumerate}

\dmsection{Purpose}

This memorandum documents two design decisions for the PCFPlan increment editor,
arising from issue \#30 (GitLab \#20, reported by K. Larson): (a) how the four
displayed increments are selected, and (b) how a payload-configuration copy
between increments resolves collisions. Both are captured as proposed
requirements pending baseline: REQ-011a (issue \#32) and REQ-060 (issue \#33).

\dmsection{Background}

PCFPlan displays four ISS increments (timeline and table columns) selected by a
date-driven window --- the current increment (by today's date) plus the next
three (REQ-011, \texttt{pdf\_generator.py:\_select\_increments}). The user cannot
choose which four are shown. It was verified against the local database that
deactivating the current increment shifts the window's left edge \textit{backward}
to the previous still-active increment (an older one), which is the mechanism
behind the unexpected columns reported in \#30. Separately, a newly created
increment's payload configuration starts empty, forcing manual rebuild; the
reporter requested the ability to copy a prior increment as a starting point.

\dmsection{Decision 1 --- Increment Display-Window Selection (REQ-011a, \#32)}

\textbf{Options considered:}
\begin{itemize}
    \item \textbf{1A --- Starting-increment selector (SELECTED).} The user picks a
    starting increment; the four displayed are that increment plus the next three.
    \item \textbf{1B --- As-of date override.} The user sets an effective date; the
    window stays date-driven but anchored to that date.
    \item \textbf{1C --- Pin four increments individually.} The user selects each of
    the four columns explicitly.
\end{itemize}

\textbf{Decision:} Option 1A. The user selects the starting increment; the
displayed four are that increment plus the next three. The selection is
\textbf{sticky} --- persisted in the database (per the project convention that all
editable state lives in the database, not config files) so it survives session to
session. Absent an explicit selection, the system retains the default REQ-011
behavior (current increment plus next three).

\textbf{Rationale:} 1A is the simplest mental model and matches how the user
reasoned about the problem. It reuses the existing ``anchor plus next three''
windowing with a user-chosen anchor instead of today's date, minimizing UI and
code change. 1B and 1C add complexity (date semantics; four-way selection state)
without additional value for the stated need.

\dmsection{Decision 2 --- Payload-Config Copy Collision Behavior (REQ-060, \#33)}

\textbf{Context:} Copying a source increment's items into a target increment can
collide with the items table constraint
\texttt{UNIQUE(increment, module\_key, parent\_name, item\_name)}. A collision
occurs only when the target already contains an item with the same module,
parent, and name. In the primary use case the target is empty, so no collision
arises.

\textbf{Options considered:}
\begin{itemize}
    \item \textbf{2A --- Skip duplicates.} Keep the target's existing item; ignore
    the incoming copy.
    \item \textbf{2B --- Overwrite, warn if not empty (SELECTED).} Source items
    replace colliding target items; if the target already contains items, warn the
    user (confirmation) before proceeding.
    \item \textbf{2C --- Empty-target-only.} Refuse the copy (warn) if the target
    is not empty.
\end{itemize}

\textbf{Decision:} Option 2B (\textbf{recommended, pending user confirmation}). On
copy, if the target increment already contains items, the system warns the user
before proceeding; on confirmation, source items overwrite colliding target
items, and non-colliding target items are unaffected. Because this behavior
directly affects the reporter's day-to-day workflow, her preference is being
solicited on issue \#30 before REQ-060 is baselined; the decision will be
finalized on her input and is explicitly held open to revision.

\textbf{Rationale:} Overwrite makes ``copy I77 $\rightarrow$ I78'' an idempotent
``make the target match the source'' operation with predictable results. The
not-empty warning guards against accidental clobbering of manual edits. Skip (2A)
was rejected because it can silently leave a half-updated target that diverges
from the source in non-obvious ways.

\dmsection{Implementation and Constraints}

\begin{itemize}
    \item Sticky window selection persisted in the database; stdlib-only; CUI-safe;
    database auto-pushed on quit per project workflow.
    \item Window selection updates both the timeline and the table columns
    (REQ-021, which requires the table columns to match the four shown).
    \item Copy carries all item attributes (type/style, flags, note references,
    arrival/descent, hierarchy).
    \item Regression tests added per REQ-QA-001 so the defect class cannot recur.
\end{itemize}

\dmsection{Verification}

\begin{itemize}
    \item \textbf{REQ-011a (Test):} selecting a starting increment renders exactly
    those four columns in timeline and table; with no selection, behavior matches
    REQ-011.
    \item \textbf{REQ-060 (Test):} copying I77 $\rightarrow$ I78 reproduces I77's
    items under I78; edits to the copy do not affect the source; collisions
    overwrite; a non-empty target triggers the warning.
    \item Both mapped in the verification matrix (requirement $\rightarrow$ file
    $\rightarrow$ line $\rightarrow$ method) per SE Phase 4.
\end{itemize}

\dmsection{References}

\begin{itemize}
    \item Issues \#30 (report), \#32 (REQ-011a), \#33 (REQ-060), \#31 (hardening).
    \item PCFPlan requirements: REQ-011, REQ-021, REQ-026, REQ-031, REQ-032, REQ-QA-001.
    \item Systems Engineering process, Phase 3 (Decision Documentation); RFC 2119.
\end{itemize}

\end{enumerate}
}

% ============================================================================
% Decision Memorandum Base Template
% ============================================================================
% This is the shared base template for all Decision Memoranda.
% DO NOT edit this file for individual decisions.
%
% Usage: Create a thin wrapper that defines variables and content, then:
%   % ============================================================================
% Decision Memorandum Base Template
% ============================================================================
% This is the shared base template for all Decision Memoranda.
% DO NOT edit this file for individual decisions.
%
% Usage: Create a thin wrapper that defines variables and content, then:
%   % ============================================================================
% Decision Memorandum Base Template
% ============================================================================
% This is the shared base template for all Decision Memoranda.
% DO NOT edit this file for individual decisions.
%
% Usage: Create a thin wrapper that defines variables and content, then:
%   \input{_template.tex}
%
% Required variables (define before \input):
%   \UniqueID       - e.g., DM-2026-001
%   \DocumentDate   - e.g., January 19, 2026
%   \AuthorName     - Signer name
%   \AuthorTitle    - Signer title
%   \ToField        - Distribution
%   \SubjectField   - Subject line
%   \OPRField       - Office of Primary Responsibility
%   \dmContent      - The full content (enumerate environment with \dmsection items)
%
% ============================================================================

\documentclass[11pt,letterpaper]{article}

% ============================================================================
% PACKAGES
% ============================================================================
\usepackage[margin=1in, top=1.15in, headheight=30pt]{geometry}
\usepackage{graphicx}
\usepackage{fancyhdr}
\usepackage{lastpage}
\usepackage{enumitem}
\usepackage{xcolor}
\usepackage{hyperref}
\usepackage{tabularx}

% ============================================================================
% DOCUMENT CONFIGURATION
% ============================================================================
\hypersetup{
    colorlinks=true,
    linkcolor=blue,
    urlcolor=blue,
    pdfborder={0 0 0}
}

% Remove paragraph indentation
\setlength{\parindent}{0pt}
\setlength{\parskip}{6pt}

% List formatting - match Word style
\setlist[enumerate]{leftmargin=0.5in, labelwidth=0.25in, labelsep=0.1in, align=left, itemsep=12pt, topsep=12pt}
\setlist[enumerate,1]{label=\arabic*.}
\setlist[itemize]{leftmargin=0.5in, itemsep=3pt, topsep=6pt}

% ============================================================================
% HEADER AND FOOTER CONFIGURATION
% ============================================================================
\pagestyle{fancy}
\fancyhf{}

% Header: text-only, Boeing / ISS Program identification
\fancyhead[L]{\large\textbf{Decision Memorandum}\\[-2pt]\scriptsize Boeing \textbar{} International Space Station Program}
\fancyhead[R]{\small \UniqueID\\[-2pt]\scriptsize \DocumentDate}
\renewcommand{\headrulewidth}{0.4pt}

% Footer: Unique ID (left), Page X of Y (center), Date (right)
\fancyfoot[L]{\small \UniqueID}
\fancyfoot[C]{\small Page \thepage\ of \pageref{LastPage}}
\fancyfoot[R]{\small \DocumentDate}
\renewcommand{\footrulewidth}{0.4pt}

% ============================================================================
% CUSTOM COMMANDS
% ============================================================================
% Bold section headers for numbered sections
\newcommand{\dmsection}[1]{\item \textbf{#1}\par}

% ============================================================================
% BEGIN DOCUMENT
% ============================================================================
\begin{document}

% ----------------------------------------------------------------------------
% UNIQUE ID (top right)
% ----------------------------------------------------------------------------
\hfill \UniqueID

\vspace{0.25in}

% ----------------------------------------------------------------------------
% SIGNATURE BLOCK
% ----------------------------------------------------------------------------
\underline{\hspace{2.5in}}\\[2pt]
\AuthorName\\
\AuthorTitle

\vspace{0.25in}

% ----------------------------------------------------------------------------
% MEMO HEADER FIELDS
% ----------------------------------------------------------------------------
\textbf{DATE:} \DocumentDate

\textbf{TO:} \ToField

\textbf{SUBJECT:} \SubjectField

\textbf{OFFICE OF PRIMARY RESPONSIBILITY:} \OPRField

\vspace{0.25in}

% ============================================================================
% MAIN CONTENT - NUMBERED SECTIONS
% ============================================================================
\dmContent

\end{document}

%
% Required variables (define before \input):
%   \UniqueID       - e.g., DM-2026-001
%   \DocumentDate   - e.g., January 19, 2026
%   \AuthorName     - Signer name
%   \AuthorTitle    - Signer title
%   \ToField        - Distribution
%   \SubjectField   - Subject line
%   \OPRField       - Office of Primary Responsibility
%   \dmContent      - The full content (enumerate environment with \dmsection items)
%
% ============================================================================

\documentclass[11pt,letterpaper]{article}

% ============================================================================
% PACKAGES
% ============================================================================
\usepackage[margin=1in, top=1.15in, headheight=30pt]{geometry}
\usepackage{graphicx}
\usepackage{fancyhdr}
\usepackage{lastpage}
\usepackage{enumitem}
\usepackage{xcolor}
\usepackage{hyperref}
\usepackage{tabularx}

% ============================================================================
% DOCUMENT CONFIGURATION
% ============================================================================
\hypersetup{
    colorlinks=true,
    linkcolor=blue,
    urlcolor=blue,
    pdfborder={0 0 0}
}

% Remove paragraph indentation
\setlength{\parindent}{0pt}
\setlength{\parskip}{6pt}

% List formatting - match Word style
\setlist[enumerate]{leftmargin=0.5in, labelwidth=0.25in, labelsep=0.1in, align=left, itemsep=12pt, topsep=12pt}
\setlist[enumerate,1]{label=\arabic*.}
\setlist[itemize]{leftmargin=0.5in, itemsep=3pt, topsep=6pt}

% ============================================================================
% HEADER AND FOOTER CONFIGURATION
% ============================================================================
\pagestyle{fancy}
\fancyhf{}

% Header: text-only, Boeing / ISS Program identification
\fancyhead[L]{\large\textbf{Decision Memorandum}\\[-2pt]\scriptsize Boeing \textbar{} International Space Station Program}
\fancyhead[R]{\small \UniqueID\\[-2pt]\scriptsize \DocumentDate}
\renewcommand{\headrulewidth}{0.4pt}

% Footer: Unique ID (left), Page X of Y (center), Date (right)
\fancyfoot[L]{\small \UniqueID}
\fancyfoot[C]{\small Page \thepage\ of \pageref{LastPage}}
\fancyfoot[R]{\small \DocumentDate}
\renewcommand{\footrulewidth}{0.4pt}

% ============================================================================
% CUSTOM COMMANDS
% ============================================================================
% Bold section headers for numbered sections
\newcommand{\dmsection}[1]{\item \textbf{#1}\par}

% ============================================================================
% BEGIN DOCUMENT
% ============================================================================
\begin{document}

% ----------------------------------------------------------------------------
% UNIQUE ID (top right)
% ----------------------------------------------------------------------------
\hfill \UniqueID

\vspace{0.25in}

% ----------------------------------------------------------------------------
% SIGNATURE BLOCK
% ----------------------------------------------------------------------------
\underline{\hspace{2.5in}}\\[2pt]
\AuthorName\\
\AuthorTitle

\vspace{0.25in}

% ----------------------------------------------------------------------------
% MEMO HEADER FIELDS
% ----------------------------------------------------------------------------
\textbf{DATE:} \DocumentDate

\textbf{TO:} \ToField

\textbf{SUBJECT:} \SubjectField

\textbf{OFFICE OF PRIMARY RESPONSIBILITY:} \OPRField

\vspace{0.25in}

% ============================================================================
% MAIN CONTENT - NUMBERED SECTIONS
% ============================================================================
\dmContent

\end{document}

%
% Required variables (define before \input):
%   \UniqueID       - e.g., DM-2026-001
%   \DocumentDate   - e.g., January 19, 2026
%   \AuthorName     - Signer name
%   \AuthorTitle    - Signer title
%   \ToField        - Distribution
%   \SubjectField   - Subject line
%   \OPRField       - Office of Primary Responsibility
%   \dmContent      - The full content (enumerate environment with \dmsection items)
%
% ============================================================================

\documentclass[11pt,letterpaper]{article}

% ============================================================================
% PACKAGES
% ============================================================================
\usepackage[margin=1in, top=1.15in, headheight=30pt]{geometry}
\usepackage{graphicx}
\usepackage{fancyhdr}
\usepackage{lastpage}
\usepackage{enumitem}
\usepackage{xcolor}
\usepackage{hyperref}
\usepackage{tabularx}

% ============================================================================
% DOCUMENT CONFIGURATION
% ============================================================================
\hypersetup{
    colorlinks=true,
    linkcolor=blue,
    urlcolor=blue,
    pdfborder={0 0 0}
}

% Remove paragraph indentation
\setlength{\parindent}{0pt}
\setlength{\parskip}{6pt}

% List formatting - match Word style
\setlist[enumerate]{leftmargin=0.5in, labelwidth=0.25in, labelsep=0.1in, align=left, itemsep=12pt, topsep=12pt}
\setlist[enumerate,1]{label=\arabic*.}
\setlist[itemize]{leftmargin=0.5in, itemsep=3pt, topsep=6pt}

% ============================================================================
% HEADER AND FOOTER CONFIGURATION
% ============================================================================
\pagestyle{fancy}
\fancyhf{}

% Header: text-only, Boeing / ISS Program identification
\fancyhead[L]{\large\textbf{Decision Memorandum}\\[-2pt]\scriptsize Boeing \textbar{} International Space Station Program}
\fancyhead[R]{\small \UniqueID\\[-2pt]\scriptsize \DocumentDate}
\renewcommand{\headrulewidth}{0.4pt}

% Footer: Unique ID (left), Page X of Y (center), Date (right)
\fancyfoot[L]{\small \UniqueID}
\fancyfoot[C]{\small Page \thepage\ of \pageref{LastPage}}
\fancyfoot[R]{\small \DocumentDate}
\renewcommand{\footrulewidth}{0.4pt}

% ============================================================================
% CUSTOM COMMANDS
% ============================================================================
% Bold section headers for numbered sections
\newcommand{\dmsection}[1]{\item \textbf{#1}\par}

% ============================================================================
% BEGIN DOCUMENT
% ============================================================================
\begin{document}

% ----------------------------------------------------------------------------
% UNIQUE ID (top right)
% ----------------------------------------------------------------------------
\hfill \UniqueID

\vspace{0.25in}

% ----------------------------------------------------------------------------
% SIGNATURE BLOCK
% ----------------------------------------------------------------------------
\underline{\hspace{2.5in}}\\[2pt]
\AuthorName\\
\AuthorTitle

\vspace{0.25in}

% ----------------------------------------------------------------------------
% MEMO HEADER FIELDS
% ----------------------------------------------------------------------------
\textbf{DATE:} \DocumentDate

\textbf{TO:} \ToField

\textbf{SUBJECT:} \SubjectField

\textbf{OFFICE OF PRIMARY RESPONSIBILITY:} \OPRField

\vspace{0.25in}

% ============================================================================
% MAIN CONTENT - NUMBERED SECTIONS
% ============================================================================
\dmContent

\end{document}

